% Requires booktabs. Methodology supplement; not a completed-results claim.
\section{Exploratory source-held-out development after external evaluation}
\label{sec:v5-exploratory-sourceheldout}

Following the frozen Mureka and Saraga transfer experiments, a new, explicitly
exploratory development package combines 2,207 accepted 60-second excerpts:
1,311 Human and 896 AI recordings, from six Human sources and two generators.
This consists of the original 1,604 development recordings, 500 Mureka
recordings and 103 Saraga recordings. It is the eligible common-context subset,
not the entire collected corpus. The original 438 locked and 50 pilot
recordings remain excluded. A separate origin ledger preserves original
roles, metadata and admission flags; no upstream role is overwritten.

Mureka and Saraga outcomes have already been examined. Their reuse cannot be
presented as untouched confirmation of hypotheses motivated by those outcomes.
Source-held-out evaluation prevents the specified within-fit source overlap,
but does not remove research-level exposure or establish pretraining
independence. The preceding frozen experiments remain unchanged.

\subsection{Complete family and quantity ablation}

All 54 existing descriptors are retained: S (15), D (3), R (3), P (6),
F (15), H (6), and M (6). Every recording and its missing values remain in
the common cohort. The primary design evaluates all 127 nonempty subsets at
five caps, with training-only median imputation and missingness indicators.
Two additional all-cap diagnostic modes use median-imputed values without
indicators, or missingness indicators alone. These diagnostics do not select
features or models.

With seed 20260907, global groups are assigned to five hash buckets. Each
source-held-out arm excludes the entire held-out source and every current
test-bucket group from training. Its test set contains current-bucket examples
from the held source and the opposite class. The verified schedule contains
28 Human-source folds, 10 generator folds, and five ordinary group-disjoint
descriptive folds. Saraga occupies only three buckets, with 77, 14 and 12
Human recordings; its two empty source-fold cells are recorded, not filled.

A cap limits groups per class/source within the training fold, retaining all
recordings in selected groups. It is not a recording count or a balanced
total sample size. Deterministic group rankings are nested across caps.
Actual training-row ranges over the 43 valid folds are shown in
Table~\ref{tab:v5-caps}.

\begin{table}[htbp]
\centering
\caption{Verified training quantities in the prospective v5 schedule.
The cap counts groups per training source, not songs.}
\label{tab:v5-caps}
\begin{tabular}{lrr}
\toprule
Group cap & Minimum training rows & Maximum training rows \\
\midrule
25 & 194 & 389 \\
50 & 355 & 678 \\
100 & 578 & 898 \\
200 & 963 & 1,297 \\
All & 1,351 & 1,805 \\
\bottomrule
\end{tabular}
\end{table}

The fixed weighted ridge model has penalty 10 and an unpenalized intercept.
Training weights balance classes, sources within class, groups within source,
and recordings within group, then sum to training size $n$. Consequently,
the mean-loss regularization is $10/n$: quantity also changes effective
regularization. This design preserves the historical model implementation,
but is not a causal quantity-only or diversity-only experiment. All decisions
use raw, unbounded identity-link score $\geq 0.5$; these scores are not
calibrated probabilities. No full-cohort refit or threshold search is performed.

\subsection{Aggregation and acceptance}

Within each held-source arm, candidate, cap, mode and source, the recording-level
class-conditional correct rate is the correct count divided by unique
out-of-fold recordings: Human specificity or AI sensitivity, not joint accuracy.
The component-level endpoint is the unweighted mean of individual global-group
correct rates. Averaging fold-level component means would be incorrect when
buckets contain different numbers of components, as in Saraga's 2/1/2 split.

Two-class summaries average source-pair metrics within a fold, valid folds
within a held-source arm, and then held-source arms within a fold type.
Human-source, generator and ordinary group holdouts remain separate.
Pooled-fold companion metrics are explicitly distinguished. No pooled
out-of-fold AUC is computed across separately fitted score scales, and no
raw predictions are pooled across unlike source-held-out arms. All 635
primary subset/cap cells per reporting arm are retained; no winner or
statistical significance is inferred from descriptive fold ranges.

The frozen schedule specifies 27,305 primary and 10,922 diagnostic fits,
with 10,687,558 predictions in total. Every training transformation, model,
test identity and output metric is persisted. At the time this methodology
supplement was written, execution was in progress: neither completion nor
independent numerical acceptance is claimed. Final performance reporting
requires a complete committed publication and independent score and metric
reconstruction.

\subsection{Reproducibility references}

All paths below are relative to
\path{artifacts/audio_phenomena_expansion_20260907/}:
\begin{itemize}
\item Protocol: \path{EQUAL60_EXPLORATORY_V5_PROTOCOL_EN.md}.
\item Derived-data code: \path{code/prepare_evaluation_inputs_v5.py};
evaluation code: \path{code/evaluate_new_phenomena_v5.py}.
\item Inputs and original-role ledger:
\path{results/equal60_exploratory_package_v5_v1/}.
\item Parent input acceptance:
\path{audit/exploratory_v5_package_parent_acceptance_v1.json}.
\item Reviewed split schedule:
\path{results/equal60_exploratory_v5_draft_v1/}.
\item Execution authorization:
\path{preregistration/equal60_exploratory_v5_frozen_v1.json}.
\item In-progress outputs, not yet accepted results:
\path{results/equal60_exploratory_v5_results_v1/}.
\item Existing third-generator context constraints:
\path{THIRD_GENERATOR_CONTEXT_ELIGIBILITY_EN.md}.
\end{itemize}
